A bored cosmologist comes up with a thought experiment to determine the Hubble
constant ($H_0$) for his model of a Steady-State-Universe. In this experiment,
a large, fully reflecting flat mirror -- carrying several gyroscopes that
would maintain its spatial orientation in the same plane -- would be placed at
a distance $D$ from the Solar System in a region without gravitational
influences. From the Earth, a laser beam would be directed towards that region
for a long period of time. After a time $T$, the radiation would return and be
detected, allowing the determination of the fixed constant $H_0$.

\begin{description}
\item[(a)] (7 points) Find an expression for $H_0$ as a function of $D$, $c$
  (speed of light) and $T$. Consider that the separation $S$ between the Solar
  System and the mirror increases only due to the expansion of the universe
  according to the law $S = s\,e^{H_0 t}$, where $s$ is the initial
  separation. You may use $e^x \approx 1 + x$ for $x \ll 1$, if necessary.
\item[(b)] (3 points) Imagine that such a mirror is located in the vicinity of
  the star Vega (which also features on the logo of the 1st GeCAA). Vega was
  the first star outside the Solar System to be photographed and one of the
  first stars whose parallax ($p = 0.125''$) was accurately measured in 1840
  by G.~W.~von Struve. Estimate the total duration of this $H_0$ measurement
  experiment.
\end{description}
